\subsection{Optical/IR Spectroscopy}

\newtext{With a reliable host identified, we observed each galaxy with optical/IR spectroscopic instruments. Detection of line emission allows} \oldtext{Optical spectroscopy is crucial for} accurate redshift measurement, line-flux estimates to better understand the nature of ongoing star-formation and nuclear activity, and for modeling the ages of the stellar populations in a galaxy. The age of the stellar population of an FRB host galaxy \newtext{tells us about how star formation creates FRB sources, a crucial test of origin models.}\oldtext{, and preferably host environment, encodes the age of the probable progenitor object, which can be determined by probing host-galaxy star-formation histories.}
Table \ref{tab:details} summarizes the spectroscopic instruments used for this FRB sample.

We observed with the Low-Resolution Imaging Spectrometer on the Keck I telescope~\citep[Keck-I/LRIS][]{1995PASP..107..375O} on 26 May, 21 July, and 17 October of 2022. \oldtext{On the last of these observing sessions, the red component of the detector had malfunctioned. We used a mirror to direct light into the blue arm, and} The light was dispersed using a 300/5000 grism to measure a detector-frame wavelength range from 3450 to 6700\AA. The spectra were reduced with the standard \sw{lpipe} software~\citep{2019PASP..131h4503P} and calibrated using observations of the standard star BD+28 4211. The spectrum was scaled to match the PS1 photometry to account for slit losses. 

We observed with the Double Spectrograph~\citep[DBSP;][]{1982PASP...94..586O} mounted at the Palomar Hale 200-inch telescope on 1 June, 12 December of 2022 and 28 May 2023. Data were reduced following procedures described in \citet{2022MNRAS.513..982R}. 

We measure the spectroscopic redshift and the line fluxes of the host galaxies using the penalized PiXel-Fitting software~\citep[\sw{pPXF};][]{2017MNRAS.466..798C, 2022arXiv220814974C} by jointly fitting the stellar continuum and nebular emission using the MILES stellar library~\citep{2006MNRAS.371..703S}. \newtext{The line fluxes are shown in Table \ref{tab:lines}. A redshift could be reliably identified for all hosts by the detection of at least 3 emission lines. For one FRB, \frbjackie, we measured the redshift by modeling absorption lines.}
